\section{Trayectorias de búsqueda y modelado potenciado por búsqueda}

Yuandong Tian continúa desarrollando esta línea central; el primer paso es tratar el propio proceso de búsqueda como una señal de supervisión, lo cual conserva la capacidad generativa del modelo neuronal y a la vez le da a los resultados una estructura interpretable.

\subsection{Trayectorias de búsqueda como weak supervision}

El expositor mostró la curva de entrenamiento del modelo potenciado por búsqueda: dentro del mismo orden de magnitud, el modelo de trayectorias de búsqueda alcanza un rendimiento superior al 80\%, mientras que el modelo puramente generativo se mantiene alrededor del 20\%. Esto se debe a que los tokens de salida del primero no solo contienen el plan final, sino también información como el elaborated trace y evaluaciones de cost, lo cual le brinda al Transformer más contexto.
\begin{figure}[H]
\centering
\includegraphics[width=\textwidth]{frame-search-curve.jpg}
\caption{Training curve del modelo potenciado por búsqueda comparado con el modelo puramente generativo, contrastando el baseline de 175M parámetros con el trace model de 15M parámetros. \protect\footnotemark}
\end{figure}
\footnotetext{Intervalo de tiempo en el video: 00:25:05--00:25:18.}
\begin{knowledgebox}{El valor de las trayectorias de búsqueda}
Las trayectorias de búsqueda permiten que el modelo «piense» antes de generar cada paso: puede reproducir la ruta de ejecución, cuantificar heuristics e incluso agregar restricciones estáticas, lo cual reduce la tasa de error durante la inference y hace viable la revisión visual.
\end{knowledgebox}

El expositor también enfatizó que, en los experimentos reales (Soang, Maze, entre otros), los datos de entrenamiento no son instrucciones vacías de un solo paso, sino el trace, los valores de costo y el plan generados por el solver. Al muestrear 64 planes sobre el trace y luego reordenarlos con heuristics, el modelo aprende cadenas de razonamiento más largas con pocas muestras.

\subsection{Search Former y DU Forer}

Sobre esta base, el equipo propuso Search Former e introdujo el trace drop en DU Forer. Al eliminar aleatoriamente parte de los tokens de búsqueda (por ejemplo, ciertas estimaciones de cost o intermediate node), el modelo aprende a recuperar el núcleo de la solución incluso cuando faltan datos, de manera similar al entrenamiento contra interferencias self-supervised.
\begin{importantbox}{La lógica de entrenamiento de DU Forer}
1) primero se usa el solver para muestrear el trace completo; 2) se hace drop de tokens específicos por nivel (por ejemplo, cost, create, exploration); 3) se hace que el Transformer reconstruya el plan a partir del trace residual; 4) se actualizan el Transformer y la drop layer mediante retropropagación del loss.
\end{importantbox}
\begin{enumerate}
    \item \textbf{Level 0:} Se conserva el trace completo como salida de línea base.
    \item \textbf{Level 1:} Se elimina el create cost y solo se conserva el path structure.
    \item \textbf{Level 2:} Se eliminan algunas acciones de exploración, lo que refuerza el modelado de los puntos de inflexión clave.
    \item \textbf{Level 3:} Solo se conserva el plan token final, y el modelo debe reconstruir por sí mismo el trace completo.
\end{enumerate}
\begin{figure}[H]
\centering
\includegraphics[width=\textwidth]{frame-dspy-architecture.jpg}
\caption{Arquitectura de DU Forer: Search Former genera el trace, que luego pasa por la drop layer y el policy head para producir el plan. \protect\footnotemark}
\end{figure}
\footnotetext{Intervalo de tiempo en el video: 00:32:25--00:32:45.}
\begin{warningbox}{Longitud del trace y equilibrio de costos}
Un trace largo aporta más contexto, pero también exige mayor longitud de tokens, memoria de GPU y sobrecarga de análisis; en la práctica es necesario controlar la estrategia de drop para evitar secuencias que sean matemáticamente completas pero inutilizables durante la inference.
\end{warningbox}

\subsection{Resumen de la sección}

Usar el trace de búsqueda como objetivo de aprendizaje, y regular la densidad de información mediante la estrategia de drop de DU Forer, permite que incluso los modelos pequeños capturen con eficiencia la calidad de la planificación.

